
\chapter{Процессинговый центр}\label{ch:processing-center}
Каждое архитектурное решение процессингового центра проверяется
двумя вопросами: что произойдёт с~задержкой и доступностью --
и~кто контролирует ключевой материал.

\section{Компоненты процессингового центра}\label{sec:pc-components}

Платёжный шлюз принимает запрос ТСП и переводит его
в~сообщение ISO~8583. Между шлюзом и карточной сетью стоит
процессинговый центр -- система, которая выбирает маршрут,
выполняет криптографические проверки, записывает след операции
и~позже превращает авторизацию в~клиринговую и расчётную запись.

\definitionnote{Процессинговый центр}{Система, которая авторизует
  карточные операции, формирует клиринг, считает расчётные позиции и
  выполняет криптографические операции. В~терминах Закона о
  НПС он часто совмещает роли операционного и платёжного клирингового
центра.}

Процессинг состоит из четырёх узлов: коммутатор авторизации,
криптографический модуль, клиринговый модуль и расчётный модуль.

\subsection*{Коммутатор авторизации}

Коммутатор авторизации принимает каждое входящее
сообщение ISO~8583 от эквайеров или от шлюзов, действующих
от их имени, выбирает маршрут к~эмитенту или в~карточную сеть
и возвращает ответ. Работает поштучно и~в~реальном времени:
допустимая задержка измеряется не секундами,
а десятками миллисекунд.~\cite{visa-core-rules,mc-tpr}

Маршрут внутри коммутатора определяется BIN-таблицей:
для каждого BIN-диапазона известно, через какой канал отправлять
запрос -- VisaNet, Banknet, ОПКЦ НСПК или прямое подключение
к~эмитенту. Карточные сети регулярно публикуют изменения
к BIN-таблице, и процессинг обязан применять их в~установленные сроки.
Иначе часть транзакций уйдёт не туда и~будет отклоняться без вины
клиента или ТСП.~\cite{visa-core-rules,mc-rules}

Коммутатор извлекает BIN -- первые 6--8 цифр PAN из DE~2 --
и ищет совпадение в BIN-таблице, определяя карточную сеть:
Visa, Mastercard или Мир. Далее проверяет, принадлежит ли BIN
тому же банку, в~котором работает процессинг. Такой случай
называется on-us: сообщение маршрутизируется напрямую
к~авторизационному хосту эмитента, минуя сеть, а~interchange
не начисляется. В~режиме off-us коммутатор выбирает исходящий
канал -- VisaNet, Banknet или ОПКЦ, -- форматирует сообщение
по Interface Specification конкретной сети и отправляет запрос.
Различия между сетями проявляются в DE~25, DE~48 и~subelement
layout. Ответ возвращается тем же маршрутом.

Помимо маршрутизации, коммутатор выполняет базовую валидацию
сообщения: проверяет обязательные поля, формат MTI
и корректность bitmap, дополняет запрос данными, которые
требует конкретная сеть, и~ведёт журнал транзакций.
Запись с~точностью до миллисекунд нужна для отладки
и~для последующей сверки с~клиринговыми файлами.

Когда внешний маршрут недоступен, включается резервная
авторизация \enquote{stand-in processing}. Чаще всего \textit{stand-in}
выполняет сама карточная сеть: VisaNet и Mastercard
предоставляют STIP-сервис, который отвечает за эмитента,
если тот недоступен из-за таймаута или сетевого сбоя.
Процессинговый центр на стороне эмитента тоже может реализовать
локальный \textit{stand-in}, опираясь на собственные правила: лимит
суммы, BIN-диапазон, историю операций по карте.
В~обоих случаях \textit{stand-in} временно берёт на себя кредитный риск,
который в~обычном режиме несёт эмитент. Поэтому его ограничивают
малыми суммами и~заранее определёнными категориями
транзакций.~\cite{visa-core-rules,mc-rules}

Ответственность за решение в~режиме \textit{stand-in} несёт тот, кто его принял.
Если \textit{stand-in} выполняла карточная сеть (VisaNet STIP,
Mastercard STIP), именно сеть берёт на себя кредитный риск
по одобренным операциям в~пределах параметров \textit{stand-in}:
лимита суммы, BIN-списка, MCC-ограничений. Если \textit{stand-in}
реализован локально на стороне процессора, риск остаётся
на процессоре или эмитенте -- в~зависимости от договора.
Параметры \textit{stand-in} -- максимальная сумма, разрешённые MCC,
BIN-диапазоны -- задают предел финансовой ответственности.
Каждый параметр за пределами нормы -- это потенциальный убыток
без покрытия.

\subsection*{Пул HSM}

Устройство HSM, иерархия ключей, key ceremony и~базовые операции --
ARQC verification, ARPC generation, PIN translation -- разобраны
в~главе~\ref{ch:crypto}. Здесь -- свойства, относящиеся именно
к~процессингу.

Пул HSM -- самое жёсткое ограничение для масштабирования процессинга.
Криптографические операции нельзя разносить по узлам так же свободно,
как сервисы без состояния, а~ключевой материал загружается через
регламентированные процедуры. Процессинговые центры строят пул HSM
с~балансировкой нагрузки и~заранее планируют ёмкость этого пула.
Конкретная производительность -- число PIN-переводов, RSA-подписей
и~ARQC-проверок в~секунду -- зависит от модели, редакции HSM
и~лицензированных функций. Её сверяют по datasheet
производителя~\cite{pci-pin-security,pci-pts-hsm,ansi-x9-24}.
К~зарубежным моделям Thales payShield, Utimaco и~Atalla
для российского рынка добавляется отдельный класс сертифицированных
ФСБ России СКЗИ -- от КриптоПро HSM до решений Аладдин Р.Д.\
и~ИнфоТеКС. Без них эксплуатация криптографии по ГОСТ невозможна.
При росте нагрузки узким местом становится HSM, а~не коммутатор
или сеть.

\importantnote{%
  HSM -- единственный компонент процессинга, у~которого горизонтальное
  масштабирование ограничено физически: каждое новое устройство
  требует key ceremony (см.~\ref{sec:crypto-key-ceremony}). Планирование
  ёмкости пула учитывает не только текущую нагрузку, но и~часы,
  необходимые на ввод устройства в~эксплуатацию.
}

Синхронизация KCV (Key Check Value) между устройствами подтверждает,
что все HSM пула работают с~одним и~тем же ключом, не раскрывая
самого ключа. При отказе одного HSM ключевой материал не теряется --
он загружен во все устройства пула во время той же церемонии. Цена
такой устойчивости -- невозможность нарастить ёмкость за минуты:
церемония занимает часы.

\subsection*{Клиринговый модуль}

Клиринговый модуль, в~отличие от авторизации, работает
пакетами. Он запускается по расписанию, собирает
авторизованные транзакции за период и добавляет данные,
которых на этапе авторизации ещё не было: финальную сумму
после добавления чаевых, результат адресной проверки
и~параметры DCC-конверсии.

Затем модуль формирует клиринговый файл в~формате сети:
TC33 для Visa, IPM для Mastercard
или собственный формат НСПК (см. главу~\ref{ch:reconciliation}).~\cite{visa-core-rules,mc-tpr,nspk-mir-rules}

На этом же этапе система сопоставляет каждую клиринговую
запись с~авторизацией по полям STAN, RRN, сумме
и~дате. Так обнаруживаются расхождения: \textit{orphan authorization} --
авторизация без клиринга, \textit{late presentment} и~\textit{force post} --
клиринг без авторизации, частичное подтверждение или
\textit{tip adjustment} -- несовпадение сумм. Каждое такое расхождение --
потенциальная потеря денег, поэтому клиринговый модуль формирует
отчёт об исключениях, который операционная команда разбирает вручную
или по заранее заданным правилам.

\subsection*{Расчётный модуль}

Расчётный модуль считает итоговые нетто-позиции участников:
сколько каждый эмитент должен перечислить и сколько каждый
эквайер должен получить с~учётом interchange, комиссий сети
и корректировок вроде чарджбэков, повторных представлений и сборов.
Результатом становится набор проводок, который уходит
в~расчётную инфраструктуру: для внутрироссийских транзакций
через НСПК и платёжную систему Банка России
(см. главу~\ref{ch:banks}), а~для трансграничных сценариев
через расчётные механизмы, предусмотренные правилами
соответствующей
схемы.~\cite{cbr-ps,visa-core-rules,mc-switching-services}

Ошибка на расчётном этапе -- прямая финансовая потеря:
именно здесь деньги начинают двигаться между банками.

\section{Производительность и отказоустойчивость}\label{sec:pc-performance}

Требования к~задержке и доступности процессингового центра
задаются внешним контрактом сети. Visa публикует maximum time limits
for Authorization Request Response и связывает просрочку
ответа со включением Stand-In Processing; Mastercard
в Transaction Processing Rules задаёт authorization
performance standards для участников сети.~\cite{visa-core-rules,mc-tpr}

\subsection*{Бюджет задержки}

Карточные сети нормируют внешний предел времени ответа, но не дают
универсальной внутренней таблицы \enquote{сколько миллисекунд
должен занимать каждый сервис}. Внутренний бюджет обязан
оставлять запас на сетевой путь, обработку у~эмитента
и~возврат ответа. Конкретная раскладка времени на парсинг,
HSM, логирование и сериализацию зависит от реализации
и измеряется внутри собственного процессинга. Типичная раскладка
для внутреннего бюджета: парсинг и валидация сообщения
1--3~мс, обращение к HSM на ARQC или PIN-перевод 3--7~мс,
BIN-lookup и маршрутизация 1~мс, логирование 1--2~мс,
сериализация ответа 1~мс. Внутренняя обработка занимает 7--15~мс;
остальное время бюджета уходит на сетевой путь до эмитента и~обратно.

\subsection*{Резервирование и геораспределение}

Процессинговый центр не может позволить себе плановый простой:
даже несколько минут недоступности -- это тысячи отклонённых
транзакций и штрафы от карточных сетей. Правила сетей не
предписывают единственную допустимую схему -- active-active,
active-standby или гибридную, -- но в~любом варианте оператору
нужно переживать отказ площадки без потери управления трафиком.

Active-active предпочтительнее active-standby по одной причине:
при переключении с~пассивного резерва площадка проходит путь
от холодного состояния до полной нагрузки -- HSM загружает ключи,
BIN-таблицы синхронизируются, очереди восстанавливаются. Каждая
секунда этого пути -- отказанные транзакции. В active-active оба
узла всегда работают под реальной нагрузкой, и~при отказе одного
трафик перераспределяется без времени раскрутки. Цена -- риск
split-brain при разрыве репликации. Для авторизации, почти не
имеющей состояния между запросами, последствия кратковременного
split-brain ограничены. Для клиринга и~расчёта нужна более жёсткая
гарантия согласованности: авторизационный след можно досверить
по внешним журналам сети, а потерянная клиринговая или расчётная
запись бьёт по денежному движению.

Основная сложность -- синхронизация состояния между ЦОД.
Журнал транзакций, BIN-таблицы, лимиты для \textit{stand-in} и HSM-ключи
обязаны оставаться согласованными в~обоих центрах.

\importantnote{%
  Авторизация и клиринг/расчёт требуют разных целей
  восстановления: первой важнее быстро восстановиться, второму важнее
  не потерять финансово значимую запись. Поэтому их обычно реплицируют по-разному.

}
\practicenote{%
  Шардинг по BIN-диапазонам остаётся простейшей и наиболее
  распространённой стратегией горизонтального масштабирования
  коммутатора авторизации. BIN однозначно определяет эмитента
  и карточную сеть, а~значит, маршрут и~набор ключей;
  транзакции с~разными BIN не зависят друг от друга
  и~могут обрабатываться разными экземплярами коммутатора
  без координации. При необходимости шард можно перенести
  на другой сервер или в~другой ЦОД без влияния на остальные
  BIN-диапазоны.

}
\subsection*{Аварийное восстановление}

Георезервирование -- размещение ЦОД в~разных географических
точках -- защищает уже не от отказа отдельного сервера,
а~от региональных катастроф: пожара, наводнения, отключения
электричества в~районе. Расстояние между центрами всегда
компромиссно. Чем дальше они друг от друга, тем лучше защита
от локальных событий, но тем выше задержка синхронной репликации.
Конкретное расстояние выбирают по сочетанию регионального риска,
качества каналов связи и~цены синхронной репликации.

RTO (Recovery Time Objective) -- время восстановления
после полного отказа ЦОД. Для процессинга оно измеряется
в~коротком операционном горизонте. Переключение трафика
обязано происходить достаточно быстро, чтобы локальный сбой
не превратился в~массовый поток отказов. Конкретное значение
определяется архитектурой и обязательствами оператора.

\section{Подключение к~платёжным системам}\label{sec:pc-connectivity}

\subsection*{VisaNet}

Подключение к VisaNet, глобальной сети Visa, определяется
правилами Visa и региональными эксплуатационными регламентами.
В~публичных документах Visa отдельно описаны использование
VisaNet, требования к~сетевым сертификатам и правила обработки
транзакций. Архитектурно VisaNet разделяет авторизацию в~BASE~I
и пакетный клиринг в~BASE~II.~\cite{visa-core-rules}

Подключение включает канал связи и сертификацию под конкретный
интерфейс сети. Публичные правила Visa оговаривают, что часть деталей,
связанных с~безопасностью сети и проприетарной внутренней
обработкой,
в~публичной редакции не раскрывается.~\cite{visa-core-rules}

\subsection*{Banknet}

Mastercard использует собственную сеть Banknet.
Публичные материалы Mastercard описывают
коммутационный уровень как отдельную систему авторизации,
клиринга и~расчёта, а Transaction Processing Rules
одновременно оперируют Authorization Request/0100,
Financial Transaction Request/0200, First Presentment/1240
и форматами IPM.~\cite{mc-switching-services,mc-tpr}

В Mastercard-среде
сосуществуют single- и dual-message сценарии, а клиринговая
часть опирается на IPM и GCMS.~\cite{mc-switching-services,mc-tpr}

\subsection*{ОПКЦ НСПК}

Все внутрироссийские карточные транзакции упираются в~НСПК --
независимо от того, выпущена ли карта внутри российской системы
Мир или относится к~международной схеме, работающей внутри России
через российскую инфраструктуру. Банк России прямо
указывает, что внутрироссийская обработка операций международных
платёжных систем ведётся через НСПК; сама НСПК публично описывает
себя как оператора ПС Мир и~ОПКЦ НСПК.~\cite{cbr-ps,nspk-about-company}

Для российского банка выход во внутреннюю
карточную инфраструктуру означает подключение
к~ОПКЦ НСПК и соблюдение правил ПС Мир и связанных
технических документов НСПК.~\cite{nspk-mir-rules}

\subsection*{СБП}

Подключение к~СБП архитектурно отличается от карточного
процессинга. СБП работает как сервис платёжной системы Банка России.
Банк России публикует отдельный порядок подключения кредитных
организаций. НСПК публикует материалы по сценариям для бизнеса,
оплате по QR, кнопке, NFC
и возвратам для бизнеса.~\cite{cbr-sbp-actions,sbp-main,sbp-faq-business}

Процессинг, который поддерживает и карты, и~СБП,
совмещает в~одной инфраструктуре два разных класса потоков:
карточный обмен на сообщениях ISO~8583 и~сервис перевода средств
по правилам СБП.~\cite{cbr-sbp-actions,sbp-main}

\subsection*{Перегрузка и управляемая деградация}

Когда карточная сеть не отвечает, процессинговый центр не может
просто ждать: тысячи новых авторизаций продолжают поступать,
и~очередь быстро переполнится.

Типовая защита -- таймауты внешних вызовов и автоматическое
размыкание канала при массовых сбоях. Точные пороги задаются
конкретной реализацией процессинга.

\section{Лицензирование и сертификация}\label{sec:pc-licensing}

Процессинговый центр не может работать без лицензии
регулятора, сертификации у~карточных сетей и~аудита PCI DSS.

\subsection*{Требования Банка России}

Закон о~НПС разводит роли операционного центра (ОЦ),
платёжного клирингового центра (ПКЦ), расчётного центра (РЦ)
и~других операторов услуг платёжной инфраструктуры (ОУПИ).
Именно в~этой терминологии процессинговый центр получает
свой юридический статус. Сами ОУПИ включаются в~реестр
операторов услуг платёжной инфраструктуры, который ведёт
Банк России.~\cite{fz-161,cbr-ps-glossary}\footnote{Точные
реквизиты порядка ведения реестра ОУПИ и~состав публикуемых
сведений сверяются по актуальной редакции на сайте Банка
России перед публикацией главы.}

Базовая нормативная рамка (821-П, 851-П, ГОСТ~Р~57580.1, 152-ФЗ
о~локализации персональных данных) разобрана
в~\ref{sec:regulation-cbr-provisions}. К~процессинговому центру
дополнительно применяется отдельная норма для платёжной системы
Банка России, включая ОПКЦ СБП и~ССНП, -- Положение Банка России
от 25.07.2022 № 802-П~\cite{cbr-802p}.

Для участия в~национальной платёжной системе \enquote{Мир}
процессинговый центр проходит сертификацию процессинга
для ПС~\enquote{Мир}; её проводит НСПК как оператор
ПС~\cite{nspk-mir-rules}. Для участия в~СБП требуется
интеграция с~ОПКЦ СБП на стороне
НСПК~\cite{cbr-sbp-rules}. Эти процедуры дополняют, а~не
заменяют требования Банка России.

\subsection*{Сертификация у~карточных сетей}

Каждая карточная сеть проводит собственную программу
сертификации для процессинговых центров, подключающихся
к~её инфраструктуре. Публичные правила Visa и Mastercard
фиксируют рамки участия в~сети: правила использования VisaNet,
требования к~обработке транзакций, составу данных, клирингу
и~расчёту.~\cite{visa-core-rules,mc-rules,mc-tpr}

Сертификация охватывает корректность форматов и обязательных
элементов данных, криптографические процедуры и устойчивость
к~операционным сбоям. Ни Visa, ни Mastercard в~открытых
редакциях не раскрывают полную матрицу тестов, связанных
с~безопасностью. Детальные приёмочные сценарии хранятся
в~закрытых руководствах и становятся доступны участнику только
на этапе подключения.~\cite{visa-core-rules,mc-rules}

\subsection*{PCI DSS}

Процессинговый центр почти неизбежно попадает в CDE
(cardholder data environment). PCI SSC отдельно разводит
cardholder data и sensitive authentication data. Для PIN,
PIN block и связанных процедур управления ключами действует
самостоятельный стандарт PCI PIN
Security.~\cite{pci-scoping-guidance,pci-faq-1335,pci-pin-security}

Требования к~подтверждению соответствия зависят
от роли и категории поставщика услуг в~программе конкретного
бренда, а~не только от формальной отнесённости к Level~1. У Visa для Level~1 service providers и VisaNet
processors используется ROC by QSA. Mastercard для
регистрируемых поставщиков услуг требует ежегодного
подтверждения соответствия PCI DSS, а~для применимых
категорий -- ежеквартальных
ASV-сканов.~\cite{visa-pci-validation-best-practice,mc-service-provider-categories-pci,pci-faq-1234-asv}
